% Generated draft: requires mathematical and source review.
\documentclass[11pt]{article}
\usepackage[margin=25mm]{geometry}
\usepackage{fontspec}
\setmainfont{Latin Modern Roman}
\usepackage{amsmath,amssymb,mathtools,unicode-math}
\setmathfont{Latin Modern Math}
\usepackage{enumitem,needspace,xurl,hyperref}
\setlength{\parindent}{0pt}
\setlength{\parskip}{5pt}
\setlength{\emergencystretch}{3em}
\newcommand{\N}{\mathbb{N}}
\newcommand{\Z}{\mathbb{Z}}
\newcommand{\Q}{\mathbb{Q}}
\newcommand{\R}{\mathbb{R}}
\newcommand{\C}{\mathbb{C}}
\newcommand{\F}{\mathbb{F}}
\newcommand{\A}{\mathbb{A}}
\newcommand{\PP}{\mathbb{P}}
\newcommand{\E}{\mathbb{E}}
\newcommand{\eps}{\varepsilon}
\newcommand{\dd}{\,\mathrm{d}}
\begin{document}
\section*{UnsolvedMath ID: 6000013}
\textbf{Source problem code:} catalogue entry
\subsection{Definitions and mathematical statement}
The source considers a closed curve on a Riemann surface carrying a projectively flat connection and asks (i) for estimates of its number of affine vertices and (ii), when its inflection points are isolated, for a relation between the vertex and inflection counts.

A conventional real-geometric interpretation would take a smooth immersion \(\gamma:S^1\to\Sigma\) and a torsion-free connection \(\nabla\) on the real tangent bundle of the Riemann surface \(\Sigma\), projectively flat in the sense that its unparametrized geodesics are locally straight in suitable coordinates. An inflection can then be described by
\[
 \dot\gamma(t)\wedge\nabla_{\dot\gamma(t)}\dot\gamma(t)=0.
\]
This description is independent of regular reparametrization. The excerpt does not, however, specify whether this real connection or a holomorphic projective-connection convention is intended.

The term ``affine vertex'' also needs its source definition. In equiaffine curve geometry it is often related to a critical point of affine curvature, defined away from inflections using an appropriate affine arclength normalization; a projectively flat connection alone does not supply all those stated conventions. No metric curvature extremum or osculating-conic definition is silently substituted here.

Once the curve class and vertex definition are fixed, write \(V(\gamma)\) for its affine-vertex count and \(I(\gamma)\) for its isolated-inflection count. Determine the intended bounds for \(V(\gamma)\) and the relation between \(V(\gamma)\) and \(I(\gamma)\). The source gives no numerical conjecture, multiplicity convention, finiteness assumption on vertices, or list of geometric/topological parameters on which the estimate may depend. In particular the isolated-inflection assumption is retained for the second request rather than strengthened to a nondegeneracy assumption on all special points.

The supplied summary reports four-vertex-type and Morse/Pl\"ucker-style relations in special convex, planar affine, and projective curve classes, but gives no precise bounds or citations. It says a sharp unified result for the stated general setting was not verified. These are qualified progress attributions, not an independently checked general theorem or current status.
\subsection{Short English statement}
Estimate affine vertices of closed curves on projectively flat Riemann surfaces, and relate their number to isolated inflections. The exact connection convention, definition of affine vertex, counting rule, and intended quantitative bounds are missing from the supplied excerpt.
\subsection{Sources}
\begin{itemize}
\item Local registry: \url{data/unsolvedmath/problems.json}, permanent UnsolvedMath ID \(6000013\), source code catalogue entry; supplied statement, background, and research summary, updated 2026-08-06. Attributed to Takeshi Sasaki, 1998.
\item Furuhata, Matsuzoe, and Urakawa, \emph{Open Problems in Affine Differential Geometry and Related Topics} (1998), item 4(b), scanned PDF pages 1--2: \url{https://www.jstage.jst.go.jp/article/iis/4/2/4_2_125/_pdf/-char/en}. Supplied manual OCR from page images, accessed 2026-07-29.
\item The summary mentions four-vertex-type and Morse/Pl\"ucker-style results without separate original-paper references. The supplied J-STAGE rights note marks the article CC BY 4.0.
\end{itemize}
\end{document}
